% GENERATED FILE. Edit canonical lessons, then run npm run generate:books.
% canonical-source-hash: fnv1a64:c95ea8161b86cb92
% canonical-lessons: IT-C441-bivio, IT-C441-vicolo, IT-C441-traversa, IT-C441-viale, IT-C441-senso-unico

\chapter{Vie e vicoli}
\label{ch:it-vie-e-vicoli}
\hlchaptermodality{441}

\begin{chapteropening}
\textbf{By the end of this chapter:} \emph{I can describe streets: a fork in the road, an alley, a side street, an avenue and a one-way street.}

Complete the last lesson of chapter 441: I can describe streets: a fork in the road, an alley, a side street, an avenue and a one-way street.
\end{chapteropening}

\section[il bivio]{il bivio — the fork (in a road)}

\label{lesson:IT-C441-bivio}



\begin{quote}
Before the new one: say the Italian for to carry on, then the Italian for to take.
\end{quote}

\subsection*{You'll want to know: il bivio}
\textbf{il bivio} — \textquotedblleft{}the fork (in a road)\textquotedblright{}.

\begin{quote}
Al bivio prendi la strada per il mare.
\end{quote}

\begin{quote}
\emph{At the fork take the road to the sea.}
\end{quote}

Figuratively, a crossroads in life: \emph{essere a un bivio}.

\begin{cousinweb}
\emph{bivio} is a learned borrowing from Latin \emph{bivium} 'a place where two roads meet', from \emph{bi-} 'two' and \emph{via} 'way' (Lewis \& Short).
\end{cousinweb}

\begin{grammarlens}[title={What you've built}]
One more word for telling, explaining and describing at B1.
\end{grammarlens}

\subsection*{Your turn}
\begin{itemize}
\raggedright
  \item \emph{Say it:} \emph{il bivio}
  \item \emph{Say it:} \emph{il bivio}, once more
  \item \emph{Say it:} a sentence of your own with \emph{il bivio}
  \item \emph{Recall:} say the Italian for to carry on, then the Italian for to take, then say \emph{il bivio} again
\end{itemize}

\begin{tcolorbox}[breakable,colback=teal!4,colframe=teal!35!black,title={Before you move on}]
What does il bivio mean? (\textquotedblleft{}The fork (in a road)\textquotedblright{}.) Say it once more.
\end{tcolorbox}
\section[il vicolo]{il vicolo — the alley, the lane}

\label{lesson:IT-C441-vicolo}



\begin{quote}
Before the new one: say the Italian for to take, then the Italian for the fork.
\end{quote}

\subsection*{You'll want to know: il vicolo}
\textbf{il vicolo} — \textquotedblleft{}the alley, the lane\textquotedblright{}.

\begin{quote}
Il ristorante è in un vicolo dietro la piazza.
\end{quote}

\begin{quote}
\emph{The restaurant is in an alley behind the square.}
\end{quote}

\emph{Un vicolo cieco} — a dead end.

\begin{cousinweb}
\emph{vicolo} is inherited from Latin \emph{vīculus} (EtymDB), a little \emph{vīcus} 'village, row of houses, street'.
\end{cousinweb}

\begin{grammarlens}[title={What you've built}]
One more word for telling, explaining and describing at B1.
\end{grammarlens}

\subsection*{Your turn}
\begin{itemize}
\raggedright
  \item \emph{Say it:} \emph{il vicolo}
  \item \emph{Say it:} \emph{il vicolo}, once more
  \item \emph{Say it:} a sentence of your own with \emph{il vicolo}
  \item \emph{Recall:} say the Italian for to take, then the Italian for the fork, then say \emph{il vicolo} again
\end{itemize}

\begin{tcolorbox}[breakable,colback=teal!4,colframe=teal!35!black,title={Before you move on}]
What does il vicolo mean? (\textquotedblleft{}The alley, the lane\textquotedblright{}.) Say it once more.
\end{tcolorbox}
\section[la traversa]{la traversa — the side street; the crossbar}

\label{lesson:IT-C441-traversa}



\begin{quote}
Before the new one: say the Italian for the fork, then the Italian for the alley.
\end{quote}

\subsection*{You'll want to know: la traversa}
\textbf{la traversa} — \textquotedblleft{}the side street; the crossbar\textquotedblright{}.

\begin{quote}
La farmacia è nella seconda traversa a destra.
\end{quote}

\begin{quote}
\emph{The chemist's is in the second side street on the right.}
\end{quote}

A street that crosses a main one. In football, \emph{la traversa} is the crossbar.

\begin{cousinweb}
\emph{traversa} is inherited from Latin \emph{trānsversa} 'lying across', from \emph{trānsvertere}, built on \emph{vertere} 'to turn' (Lewis \& Short).
\end{cousinweb}

\begin{grammarlens}[title={What you've built}]
One more word for telling, explaining and describing at B1.
\end{grammarlens}

\subsection*{Your turn}
\begin{itemize}
\raggedright
  \item \emph{Say it:} \emph{la traversa}
  \item \emph{Say it:} \emph{la traversa}, once more
  \item \emph{Say it:} a sentence of your own with \emph{la traversa}
  \item \emph{Recall:} say the Italian for the fork, then the Italian for the alley, then say \emph{la traversa} again
\end{itemize}

\begin{tcolorbox}[breakable,colback=teal!4,colframe=teal!35!black,title={Before you move on}]
What does la traversa mean? (\textquotedblleft{}The side street; the crossbar\textquotedblright{}.) Say it once more.
\end{tcolorbox}
\section[il viale]{il viale — the avenue, the boulevard}

\label{lesson:IT-C441-viale}



\begin{quote}
Before the new one: say the Italian for the alley, then the Italian for the side street.
\end{quote}

\subsection*{You'll want to know: il viale}
\textbf{il viale} — \textquotedblleft{}the avenue, the boulevard\textquotedblright{}.

\begin{quote}
Passeggiamo sul viale sotto gli alberi.
\end{quote}

\begin{quote}
\emph{We stroll along the avenue under the trees.}
\end{quote}

Wide and lined with trees.

\begin{cousinweb}
\emph{viale} is built on \emph{via} 'road', which is inherited from Latin \emph{via} 'way' (Lewis \& Short).
\end{cousinweb}

\begin{grammarlens}[title={What you've built}]
One more word for telling, explaining and describing at B1.
\end{grammarlens}

\subsection*{Your turn}
\begin{itemize}
\raggedright
  \item \emph{Say it:} \emph{il viale}
  \item \emph{Say it:} \emph{il viale}, once more
  \item \emph{Say it:} a sentence of your own with \emph{il viale}
  \item \emph{Recall:} say the Italian for the alley, then the Italian for the side street, then say \emph{il viale} again
\end{itemize}

\begin{tcolorbox}[breakable,colback=teal!4,colframe=teal!35!black,title={Before you move on}]
What does il viale mean? (\textquotedblleft{}The avenue, the boulevard\textquotedblright{}.) Say it once more.
\end{tcolorbox}
\section[il senso unico]{il senso unico — the one-way street}

\label{lesson:IT-C441-senso-unico}



\begin{quote}
Before the new one: say the Italian for the side street, then the Italian for the avenue.
\end{quote}

\subsection*{You'll want to know: il senso unico}
\textbf{il senso unico} — \textquotedblleft{}the one-way street\textquotedblright{}.

\begin{quote}
Non puoi entrare, è senso unico.
\end{quote}

\begin{quote}
\emph{You cannot go in, it is one-way.}
\end{quote}

\emph{Il senso vietato} — no entry.

\begin{cousinweb}
\emph{senso unico} is built from \emph{senso} 'sense, direction', inherited from Latin \emph{sēnsus}, and \emph{unico} 'only', inherited from Latin \emph{ūnicus}, from \emph{ūnus} 'one' (EtymDB).
\end{cousinweb}

\begin{grammarlens}[title={What you've built}]
One more word for telling, explaining and describing at B1.
\end{grammarlens}

\subsection*{Your turn}
\begin{itemize}
\raggedright
  \item \emph{Say it:} \emph{il senso unico}
  \item \emph{Say it:} \emph{il senso unico}, once more
  \item \emph{Say it:} a sentence of your own with \emph{il senso unico}
  \item \emph{Recall:} say the Italian for the side street, then the Italian for the avenue, then say \emph{il senso unico} again
\end{itemize}

\begin{tcolorbox}[breakable,colback=teal!4,colframe=teal!35!black,title={Before you move on}]
What does il senso unico mean? (\textquotedblleft{}The one-way street\textquotedblright{}.) Say it once more.
\end{tcolorbox}
